\section{基本公式}\label{節:基本公式の逆}

原始関数を探す出発点は,微分の公式を逆向きに読むことである.本節では,基本的な関数の原始関数をそうして求め,表にまとめる.以後の節では,与えられた関数をこの表のどれかに帰着させることを目標とする.

$F'(x)=f(x)$ ならば,\,$F$ は $f$ の原始関数である.この関係を使うとき,最初から正しい係数まで思い付く必要はない.形の近い候補を微分し,元の関数との差を直せばよい.

$x^r$ で試してみよう.指数を1つ増やした $x^{r+1}$ を微分すると
\[
 \frac{\dd}{\dd x}x^{r+1}=(r+1)x^r
\]
となり,余分な係数 $r+1$ が出る.そこで $r\ne-1$ ならば,両辺を $r+1$ で割って
\[
 \frac{\dd}{\dd x}\frac{x^{r+1}}{r+1}=x^r
\]
を得る.すなわち
\[
 \int x^r\dd x=\frac{x^{r+1}}{r+1}+C
\]
である.指数を1つ増やし,増やした指数で割るのは,微分で指数が1つ減り,その指数が係数として出ることを元に戻しているのである.

ここで $r=-1$ だけは別に考える必要がある.上の式へ $r=-1$ を代入すると分母が $0$ になり,意味を持たない.また,\,$x^0=1$ を微分しても $x^{-1}$ は現れない.代わりに,対数の微分
\[
 \frac{\dd}{\dd x}\log x=\frac1x\qquad(x>0)
\]
を使う.\,$x<0$ では $\log(-x)$ を微分すれば
\[
 \frac{\dd}{\dd x}\log(-x)
 =\frac1{-x}\,(-1)=\frac1x
\]
となるから,2つの形は $\log|x|$ とまとめられる.

ほかの関数も,微分の公式の左右を入れ換えて読めばよい.たとえば $a^x$ を微分すると $\log a$ が掛かるので,積分では $\log a$ で割る.\footnote{$a=1$ を除くのは,\,$\log1=0$ で割ることになるためである.\,$1^x=1$ の積分そのものは $x+C$ である.}以上をまとめると,次の表を得る.表の右の式を微分すれば,左の式に戻ることを確かめられる.

\begin{theorem}{基本的な不定積分}
次の公式は,右辺が微分でき,左辺の関数が定義される区間で成り立つ.
\[
\begin{array}{c|c}
 f(x)&\int f(x)\dd x\\ \hline
 x^r\quad(r\ne-1)&\frac{x^{r+1}}{r+1}+C\\[5pt]
 \frac1x&\log|x|+C\\[5pt]
 \e^x&\e^x+C\\[3pt]
 a^x\quad(a>0,\ a\ne1)&\frac{a^x}{\log a}+C\\[5pt]
 \sin x&-\cos x+C\\[3pt]
 \cos x&\sin x+C\\[3pt]
 \frac1{\cos^2x}&\tan x+C
\end{array}
\]
実数 $r$ に対するべきの公式は,まず $x>0$ で考える.整数のべきなど,負の $x$ でも定義される場合は,その定義域内の区間にも適用できる.
\end{theorem}

また,微分は和と定数倍を保つ.したがって,和は項ごとに積分でき,定数倍は積分の外に出せる.

ここで1つ断りをいれておこう.不定積分の $C$ は,考えている1つの区間の上での任意定数である.\,$\frac1x$ の定義域は $0$ の左右に分かれ,\,$0$ を越えて1つの区間にすることはできない.そのため,左側で加える定数と右側で加える定数が同じである必要はない.図\ref{図:対数の区間}に,この区間の分かれ方を数直線で示した.本章では以後,不定積分の等式は被積分関数が定義される各区間の上で成り立つものとし,任意定数は区間ごとに独立に選べるものとする.どの点を除き,どの区間で考えるかは,例題の解法で述べる.

\begin{figure}[htbp]
 \centering
 \begin{tikzpicture}[x=1cm,y=1cm]
  \draw[-{Stealth}] (-4,0)--(4,0) node[right] {$x$};
  \draw[very thick,blue] (-3.7,0)--(-0.12,0);
  \draw[very thick,blue] (0.12,0)--(3.7,0);
  \filldraw[fill=white] (0,0) circle (2pt);
  \node[below] at (0,-0.08) {$0$};
  \node[above] at (-2,0.18) {$\log(-x)+C_-$};
  \node[above] at (2,0.18) {$\log x+C_+$};
  \node[below] at (-2,-0.35) {$x<0$};
  \node[below] at (2,-0.35) {$x>0$};
 \end{tikzpicture}
 \caption{$0$ を越えて1つの区間にはできない.定数は左右で独立に選べる.}
 \label{図:対数の区間}
\end{figure}

\begin{bookpoint}
\textbf{POINT}\par
根号や分数を $x^r$ の形に書き直し,和と定数倍に分けて,項ごとに表の公式を使う.
\end{bookpoint}

\begin{example}{べきの積分}
$x>0$ において,\,$\int\left(2\sqrt{x}+\frac3{x^2}\right)\dd x$ を求めよ.
\end{example}

\paragraphbf{方針}
$\sqrt{x}=x^{\frac12}$,\,$\frac1{x^2}=x^{-2}$ と書き直し,べきの公式に帰着させる.

\paragraphbf{解答}
$\frac43x^{\frac32}-\frac3x+C$

\paragraphbf{解法}
$\sqrt{x}=x^{\frac12}$,\,$\frac1{x^2}=x^{-2}$ と書き,項ごとに積分すると
\begin{align*}
 \int\left(2\sqrt{x}+\frac3{x^2}\right)\dd x
 &=2\int x^{\frac12}\dd x+3\int x^{-2}\dd x\\
 &=2\cdot\frac23x^{\frac32}+3\cdot\frac{x^{-1}}{-1}+C\\
 &=\frac43x^{\frac32}-\frac3x+C
\end{align*}
となる.\,$x>0$ では両方の項が定義される.分母に $x$ があっても,\,$\frac1{x^2}$ は指数が $-2$ なので,対数ではなくべきの公式を使う.答えを微分すると
\[
 \frac43\cdot\frac32x^{\frac12}+3x^{-2}=2\sqrt{x}+\frac3{x^2}
\]
となり,係数と符号の両方が元の式と一致する.

\paragraphbf{積の積分}
和は項ごとに積分できたが,積について
\[
 \int f(x)g(x)\dd x
 =\left(\int f(x)\dd x\right)
  \left(\int g(x)\dd x\right)
\]
という規則はない.右辺を微分すると積の微分の2項が現れ,ふつうは $f(x)g(x)$ にはならない.実際,\,$f(x)=g(x)=x$ とすると,左辺は $\frac{x^3}3+C$ だが,右辺の $\frac{x^2}2\cdot\frac{x^2}2=\frac{x^4}4$ を微分すると $x^3$ になり,\,$x^2$ には戻らない.積があるときは,まず展開できるか,恒等式で和にできるかを考える.これが次節以降で扱う変形であり,後の章で学ぶ置換積分と部分積分も,この問題への別の答えである.
